% Hafta 4 — Kod sağlamlaştırmanın üç katmanı
% Derleme: py -3.12 tools/latex/derle.py docs/week-4/assets/h04-05-uc-katman.tex
\documentclass[tikz,border=8pt]{standalone}
\usepackage{cen429-cizim}
\begin{document}
\begin{tikzpicture}[node distance=8mm and 8mm]
  % Başlık
  \node[baslik] (b) at (0,0) {Kod sağlamlaştırmanın üç katmanı};
  \node[altbaslik] at (b.south west) {yalnız en alttaki katman hatayı yok eder};

  \node[uyarikutu=140mm, below=14mm of b.south west, anchor=north west] (k3)
    {\kb{3 · Kod gizleme} \textperiodcentered{} anlamayı zorlaştırır --- hatayı SİLMEZ};
  \node[kutu=140mm, below=of k3] (k2)
    {\kb{2 · Derleyici / OS korumaları} \textperiodcentered{} sömürüyü zorlaştırır --- hatayı SİLMEZ};
  \node[iyikutu=140mm, below=of k2] (k1)
    {\kb{1 · Güvenli kodlama} \textperiodcentered{} hatayı GERÇEKTEN ortadan kaldırır};

  \draw[ok] (k3) -- (k2);
  \draw[ok] (k2) -- (k1);

  % Dip şeridi
  \node[dip, text width=140mm, below=8mm of k1.south, anchor=north] {%
    Sıra: önce doğru yaz, sonra derleyiciyle sertleştir, en son gizle.};
\end{tikzpicture}
\end{document}
